\begin{lemma}[From Jensen's lemma to a superjet of each summand]
Let $E$ be a finite-dimensional real inner-product space equipped with its Borel
$\sigma$-algebra, let $N$ be a natural number and write
$V=\ell^2\bigl(\bigoplus_{j=1}^NE\bigr)$ for the $\ell^2$-product of $N$ copies
of $E$. Let $C\ge0$ and let $G_j:E\to\R$ $(1\le j\le N)$ be continuous with
$G_j(w)\le C$ for every $w$. Let $\varepsilon>0$, let $A_\varepsilon:V\to V$ be a
self-adjoint continuous linear operator, write $\mathcal G(W)=\sum_jG_j(w_j)$,
let $\mathcal G^\varepsilon$ and $G_j^\varepsilon$ be the sup-convolutions of
\noderef{SupConvolution} with parameter $\varepsilon$ formed in $V$ and in $E$,
and let $\Upsilon:V\to\R$ satisfy
\[
  \Upsilon(Z)=\mathcal G^\varepsilon(Z)
    -\tfrac12\langle A_\varepsilon Z,Z\rangle\qquad(Z\in V),
  \qquad
  \mathcal G^\varepsilon(Z)=\sum_jG_j^\varepsilon(z_j)\qquad(Z\in V).
\]
Assume $\Upsilon(0)=0$, $\Upsilon(Z)<0$ for $Z\ne0$, that $\Upsilon$ is
$\lambda$-semiconvex on $V$ in the sense of \noderef{SemiconvexOn} with
$\lambda=\varepsilon^{-1}+\|A_\varepsilon\|$, and that $\Upsilon$ is continuous
on $V$. Let $R>0$ and $\delta'>0$. Then there are $\widehat Z'\in V$, $\Pi\in V$
and symmetric continuous bilinear forms $X_j$ on $E$ $(1\le j\le N)$ with
\begin{enumerate}
\item[(J1)] $|\widehat Z'|<R$ and $|\Pi|\le\delta'+\|A_\varepsilon\|R$;
\item[(J2)] $\|X_j\|\le\varepsilon^{-1}+2\|A_\varepsilon\|$ for every $j$;
\item[(J3)] $\bigl(\Pi_j,X_j\bigr)
  \in J^{2,+}G_j\bigl(\widehat z'_j+\varepsilon\Pi_j\bigr)$ for every $j$, where
  $J^{2,+}$ is the superjet of \noderef{Superjet};
\item[(J4)] $\sum_jX_j[\zeta_j]\le\langle A_\varepsilon\zeta,\zeta\rangle$ for
  every $\zeta\in V$, where $X[h]=X(h,h)$.
\end{enumerate}
\end{lemma}
\begin{proof}
For $a\in E$ and $1\le j\le N$ let $\iota_ja\in V$ be the vector with $j$-th
component $a$ and all other components $0$; then $|\iota_ja|=|a|$ and
$\langle\xi,\iota_ja\rangle=\langle\xi_j,a\rangle$ for every $\xi\in V$, by the
definition of the $\ell^2$-product inner product (preamble). Fix an additive
Haar measure $\mu_V$ on $V$.

\emph{Step 1: a good point.}
The ball $B(0,R)$ is open and convex, and $\Upsilon$ is $\lambda$-semiconvex on
$V$, hence on the convex subset $B(0,R)$: the function
$\Upsilon+\tfrac\lambda2|\cdot|^2$ is convex on $V$ and therefore convex on
$B(0,R)$. Also $\Upsilon$ is continuous, hence continuous on
$\overline B(0,R)$, and $0$ is a strict maximum point of $\Upsilon$ on
$\overline B(0,R)$: for $Z\in\overline B(0,R)$ with $Z\ne0$ we have
$\Upsilon(Z)<0=\Upsilon(0)$. Finally $\lambda>0$ because $\varepsilon>0$ and
$\|A_\varepsilon\|\ge0$. By \noderef{JensenLemma}, applied in $V$ with this
$\lambda$, the radius $R$, the centre $0$ and the parameter $\delta'$, the set
\[
  \mathcal K=\bigl\{Z\in B(0,R):\exists\,Q\in V,\ |Q|\le\delta',\
    W\mapsto\Upsilon(W)+\langle Q,W\rangle\ \text{has a local maximum on}\
    B(0,R)\ \text{at}\ Z\bigr\}
\]
satisfies $\mu_V(\mathcal K)>0$.

By \noderef{SemiconvexSecondOrderAe}, applied in $V$ with this $\lambda\ge0$ and
the open set $B(0,R)$, there is a $\mu_V$-null set $\mathcal N$ such that every
$Z\in B(0,R)\setminus\mathcal N$ admits $q_\Upsilon\in V$ and a symmetric
continuous bilinear form $A_\Upsilon$ on $V$ with
$A_\Upsilon(\zeta,\zeta)\ge-\lambda|\zeta|^2$ for all $\zeta$ and, for every
$\delta>0$, some $\rho(\delta)>0$ with
\begin{equation}\label{eq:srjj-expansion}
  \bigl|\Upsilon(Z+\zeta)-\Upsilon(Z)-\langle q_\Upsilon,\zeta\rangle
    -\tfrac12A_\Upsilon(\zeta,\zeta)\bigr|\le\delta|\zeta|^2
  \qquad(|\zeta|<\rho(\delta)).
\end{equation}
The set $\mathcal K\setminus\mathcal N$ is nonempty: were it empty we would
have $\mathcal K\subseteq\mathcal N$, hence
$\mu_V(\mathcal K)\le\mu_V(\mathcal N)=0$ by the monotonicity of $\mu_V$,
contradicting $\mu_V(\mathcal K)>0$. So pick
$\widehat Z'\in\mathcal K\setminus\mathcal N$, a multiplier $Q$ with
$|Q|\le\delta'$ witnessing $\widehat Z'\in\mathcal K$, and expansion data
$(q_\Upsilon,A_\Upsilon)$ at $\widehat Z'$. In particular
$|\widehat Z'|<R$.

\emph{Step 2: first- and second-order conditions.}
As $B(0,R)$ is open and $\widehat Z'\in B(0,R)$, and the local maximum of
$W\mapsto\Upsilon(W)+\langle Q,W\rangle$ on $B(0,R)$ at $\widehat Z'$ holds on
some neighbourhood of $\widehat Z'$ inside $B(0,R)$, there is $\rho_1>0$ with
$B(\widehat Z',\rho_1)\subseteq B(0,R)$ and
\begin{equation}\label{eq:srjj-max}
  \Upsilon(\widehat Z'+\zeta)-\Upsilon(\widehat Z')
   \le-\langle Q,\zeta\rangle\qquad(|\zeta|<\rho_1).
\end{equation}
Let $\delta>0$ and $|\zeta|<\min(\rho_1,\rho(\delta))$. Adding
\eqref{eq:srjj-expansion} (in the form
$\langle q_\Upsilon,\zeta\rangle+\tfrac12A_\Upsilon(\zeta,\zeta)
-\delta|\zeta|^2\le\Upsilon(\widehat Z'+\zeta)-\Upsilon(\widehat Z')$) to
\eqref{eq:srjj-max} gives
\begin{equation}\label{eq:srjj-combined}
  \langle q_\Upsilon+Q,\zeta\rangle
   \le-\tfrac12A_\Upsilon(\zeta,\zeta)+\delta|\zeta|^2
   \qquad\bigl(|\zeta|<\min(\rho_1,\rho(\delta))\bigr).
\end{equation}
Fix a unit vector $v\in V$ and put $\zeta=tv$ with $t>0$ small, so that
$|\zeta|=t<\min(\rho_1,\rho(\delta))$ once $t$ is small enough. Dividing
\eqref{eq:srjj-combined} by $t$ gives
$\langle q_\Upsilon+Q,v\rangle\le t\bigl(-\tfrac12A_\Upsilon(v,v)
+\delta\bigr)$, and letting $t\downarrow0$ gives
$\langle q_\Upsilon+Q,v\rangle\le0$ for every unit $v$; applying this to $v$ and
to $-v$ gives $\langle q_\Upsilon+Q,v\rangle=0$ for every unit $v$. Writing
$u=q_\Upsilon+Q$, this says $\langle u,v\rangle=0$ for every unit $v$; if
$u\ne0$ then $v=u/|u|$ is a unit vector with $\langle u,v\rangle=|u|>0$, a
contradiction, so $u=0$. (If $V=\{0\}$ there are no unit vectors at all and the
display above is vacuous, but then $u=0$ holds outright, $V$ having only the
zero vector; the normalisation $v=u/|u|$ is legitimate precisely because it is
performed only when $u\ne0$.) Hence
\begin{equation}\label{eq:srjj-first}
  q_\Upsilon=-Q .
\end{equation}
Substituting \eqref{eq:srjj-first} into \eqref{eq:srjj-combined} leaves
$0\le-\tfrac12A_\Upsilon(\zeta,\zeta)+\delta|\zeta|^2$, i.e.\
$\tfrac12A_\Upsilon(\zeta,\zeta)\le\delta|\zeta|^2$ for
$|\zeta|<\min(\rho_1,\rho(\delta))$. Now fix $\zeta\in V$ and $\delta>0$. If
$\zeta=0$ then $A_\Upsilon(\zeta,\zeta)=0$ by bilinearity and there is nothing
to prove. If $\zeta\ne0$, put $v=\zeta/|\zeta|$, a unit vector, and apply the
displayed inequality at $tv$ for $0<t<\min(\rho_1,\rho(\delta))$: since
$A_\Upsilon(tv,tv)=t^2A_\Upsilon(v,v)$, dividing by $t^2>0$ gives
$\tfrac12A_\Upsilon(v,v)\le\delta$, and $\delta>0$ was arbitrary, so
$A_\Upsilon(v,v)\le0$; multiplying by $|\zeta|^2\ge0$ and using
$A_\Upsilon(\zeta,\zeta)=|\zeta|^2A_\Upsilon(v,v)$ gives
\begin{equation}\label{eq:srjj-second}
  A_\Upsilon(\zeta,\zeta)\le0\qquad(\zeta\in V).
\end{equation}
(The normalisation $v=\zeta/|\zeta|$ is performed only when $\zeta\ne0$, which
is what makes the argument correct when $V=\{0\}$: there every $\zeta$ is $0$
and only the first case occurs.)

\emph{Step 3: the expansion of $\mathcal G^\varepsilon$ and the forms $X_j$.}
Since $\mathcal G^\varepsilon(Z)=\Upsilon(Z)
+\tfrac12\langle A_\varepsilon Z,Z\rangle$ and, exactly, using the
self-adjointness of $A_\varepsilon$,
\[
  \tfrac12\bigl\langle A_\varepsilon(\widehat Z'+\zeta),
    \widehat Z'+\zeta\bigr\rangle
   =\tfrac12\langle A_\varepsilon\widehat Z',\widehat Z'\rangle
     +\langle A_\varepsilon\widehat Z',\zeta\rangle
     +\tfrac12\langle A_\varepsilon\zeta,\zeta\rangle,
\]
the estimate \eqref{eq:srjj-expansion} becomes, with
\[
  \Pi:=q_\Upsilon+A_\varepsilon\widehat Z'=-Q+A_\varepsilon\widehat Z',
  \qquad
  A':=A_\Upsilon+\langle A_\varepsilon\cdot,\cdot\rangle,
\]
a symmetric continuous bilinear form on $V$,
\begin{equation}\label{eq:srjj-Gexp}
  \bigl|\mathcal G^\varepsilon(\widehat Z'+\zeta)
    -\mathcal G^\varepsilon(\widehat Z')-\langle\Pi,\zeta\rangle
    -\tfrac12A'(\zeta,\zeta)\bigr|\le\delta|\zeta|^2
  \qquad(|\zeta|<\rho(\delta)).
\end{equation}
Here $|\Pi|\le|Q|+|A_\varepsilon\widehat Z'|
\le\delta'+\|A_\varepsilon\|\,|\widehat Z'|\le\delta'+\|A_\varepsilon\|R$, which
together with $|\widehat Z'|<R$ proves (J1).

Define $X_j(a,b)=A'(\iota_ja,\iota_jb)$ for $a,b\in E$; each $X_j$ is a
continuous bilinear form on $E$, and it is symmetric because $A'$ is. Taking
$\zeta=\iota_jh$ in \eqref{eq:srjj-Gexp}, and using the separability hypothesis
$\mathcal G^\varepsilon(Z)=\sum_iG_i^\varepsilon(z_i)$ --- which shows that
changing only the $j$-th component of the argument changes
$\mathcal G^\varepsilon$ only through the single summand $G_j^\varepsilon$, all
other summands being unchanged --- together with
$\langle\Pi,\iota_jh\rangle=\langle\Pi_j,h\rangle$ and $|\iota_jh|=|h|$, we get
\begin{equation}\label{eq:srjj-Gjexp}
  \bigl|G_j^\varepsilon(\widehat z'_j+h)-G_j^\varepsilon(\widehat z'_j)
   -\langle\Pi_j,h\rangle-\tfrac12X_j(h,h)\bigr|\le\delta|h|^2
  \qquad(|h|<\rho(\delta)).
\end{equation}
In particular, given $\varepsilon'>0$, apply \eqref{eq:srjj-Gjexp} with
$\delta=\varepsilon'$ and drop the absolute value on the favourable side: for
$|h|<\rho(\varepsilon')$,
\[
  G_j^\varepsilon(\widehat z'_j+h)\le G_j^\varepsilon(\widehat z'_j)
    +\langle\Pi_j,h\rangle+\tfrac12X_j(h,h)+\varepsilon'|h|^2 .
\]
Since $\varepsilon'>0$ was arbitrary and $X_j$ is symmetric, the defining
condition of \noderef{Superjet} holds, so
\begin{equation}\label{eq:srjj-jet}
  \bigl(\Pi_j,X_j\bigr)\in J^{2,+}G_j^\varepsilon(\widehat z'_j)
  \qquad(1\le j\le N).
\end{equation}

\emph{Step 4: the matrix inequality and the norm bound.}
Let $\zeta\in V$ and $\delta>0$ with $|\zeta|<\rho(\delta)$. Then
$|\zeta_j|\le|\zeta|<\rho(\delta)$ for every $j$, so summing
\eqref{eq:srjj-Gjexp} over $j$ with $h=\zeta_j$ and using the separability
hypothesis again, together with
$\sum_j\langle\Pi_j,\zeta_j\rangle=\langle\Pi,\zeta\rangle$ and
$\sum_j|\zeta_j|^2=|\zeta|^2$ (both instances of the definition of the
$\ell^2$-product inner product, preamble), gives
\[
  \Bigl|\mathcal G^\varepsilon(\widehat Z'+\zeta)
   -\mathcal G^\varepsilon(\widehat Z')-\langle\Pi,\zeta\rangle
   -\tfrac12\sum_jX_j[\zeta_j]\Bigr|\le\sum_j\delta|\zeta_j|^2=\delta|\zeta|^2 .
\]
Comparing with \eqref{eq:srjj-Gexp} and using the triangle inequality gives
$\bigl|\tfrac12A'(\zeta,\zeta)-\tfrac12\sum_jX_j[\zeta_j]\bigr|
\le2\delta|\zeta|^2$ for $|\zeta|<\rho(\delta)$. Both
$\zeta\mapsto A'(\zeta,\zeta)$ and $\zeta\mapsto\sum_jX_j[\zeta_j]$ are
homogeneous of degree two, the second because $(t\zeta)_j=t\zeta_j$ for every
$j$. Fix $\zeta\in V$ and $\delta>0$. If $\zeta=0$ then both forms vanish and
the two sides of \eqref{eq:srjj-diag} below agree. If $\zeta\ne0$, put
$v=\zeta/|\zeta|$, a unit vector, and apply the previous display at $tv$ with
$t\in(0,\rho(\delta))$; dividing by $t^2$ shows
$\bigl|\tfrac12A'(v,v)-\tfrac12\sum_jX_j[v_j]\bigr|\le2\delta$, and
$\delta>0$ was arbitrary, so
$\tfrac12A'(v,v)=\tfrac12\sum_jX_j[v_j]$. Multiplying by $|\zeta|^2$ and using
the degree-two homogeneity of both forms,
\begin{equation}\label{eq:srjj-diag}
  \sum_jX_j[\zeta_j]=A'(\zeta,\zeta)
   =A_\Upsilon(\zeta,\zeta)+\langle A_\varepsilon\zeta,\zeta\rangle
  \qquad(\zeta\in V).
\end{equation}
By \eqref{eq:srjj-second} this gives
$\sum_jX_j[\zeta_j]\le\langle A_\varepsilon\zeta,\zeta\rangle$, which is (J4);
and together with $A_\Upsilon(\zeta,\zeta)\ge-\lambda|\zeta|^2$ and the
two-sided Cauchy--Schwarz estimate
\[
  \bigl|\langle A_\varepsilon\zeta,\zeta\rangle\bigr|
   \le|A_\varepsilon\zeta|\,|\zeta|\le\|A_\varepsilon\|\,|\zeta|^2,
  \qquad\text{i.e.}\qquad
  -\|A_\varepsilon\|\,|\zeta|^2\le\langle A_\varepsilon\zeta,\zeta\rangle
   \le\|A_\varepsilon\|\,|\zeta|^2
\]
--- the lower bound being needed for the left-hand half of what follows --- it
gives
\begin{equation}\label{eq:srjj-bounds}
  -\bigl(\lambda+\|A_\varepsilon\|\bigr)|\zeta|^2
   \le\sum_jX_j[\zeta_j]\le\|A_\varepsilon\|\,|\zeta|^2\qquad(\zeta\in V).
\end{equation}
Taking $\zeta=\iota_ja$ in \eqref{eq:srjj-bounds}, so that $\zeta_j=a$ and
$\zeta_i=0$ for $i\ne j$ and $|\zeta|=|a|$, yields the bound on the diagonal
\[
  \bigl|X_j[a]\bigr|\le\bigl(\lambda+\|A_\varepsilon\|\bigr)|a|^2
  \qquad(a\in E).
\]
Each $X_j$ is a symmetric continuous bilinear form on $E$ and
$\lambda+\|A_\varepsilon\|\ge0$, since $\lambda>0$; so
\noderef{SymmetricFormNormFromDiagonal}, applied to $S=X_j$ with the constant
$C=\lambda+\|A_\varepsilon\|$, converts this diagonal bound into the same
bound on the operator norm:
\[
  \|X_j\|\le\lambda+\|A_\varepsilon\|
   =\varepsilon^{-1}+2\|A_\varepsilon\| ,
\]
which is (J2).

\emph{Step 5: transfer to $G_j$.}
The function $G_j$ is continuous on $E$ and satisfies $G_j(w)\le C\le C+|w|$,
and $\varepsilon>0$; so \noderef{SupConvolutionMagic} applies to the jet
\eqref{eq:srjj-jet} at the point $\widehat z'_j$ with the vector $\Pi_j$ and the
form $X_j$, and yields
\[
  \bigl(\Pi_j,X_j\bigr)
   \in J^{2,+}G_j\bigl(\widehat z'_j+\varepsilon\Pi_j\bigr),
\]
which is (J3).
\end{proof}
